% Canonical TeX edition. Initial text inherited from 3b067cc2eba4ac2f2f91248448727639d8ac7c5a
% Source: research/case_study/phase265_triangle_application.md
% Subsequent mathematical edits require a new audit.
% Options for packages loaded elsewhere
\PassOptionsToPackage{unicode}{hyperref}
\PassOptionsToPackage{hyphens}{url}
\documentclass[
  11pt,
  a4paper,
]{article}
\usepackage{xcolor}
\usepackage[margin=1in]{geometry}
\usepackage{amsmath,amssymb}
\setcounter{secnumdepth}{-\maxdimen} % remove section numbering
\usepackage{iftex}
\ifPDFTeX
  \usepackage[T1]{fontenc}
  \usepackage[utf8]{inputenc}
  \usepackage{textcomp} % provide euro and other symbols
\else % if luatex or xetex
  \usepackage{unicode-math} % this also loads fontspec
  \defaultfontfeatures{Scale=MatchLowercase}
  \defaultfontfeatures[\rmfamily]{Ligatures=TeX,Scale=1}
\fi
\usepackage{lmodern}
\ifPDFTeX\else
  % xetex/luatex font selection
  \setmainfont[]{STIXGeneral}
  \setmonofont[]{Menlo}
\fi
% Use upquote if available, for straight quotes in verbatim environments
\IfFileExists{upquote.sty}{\usepackage{upquote}}{}
\IfFileExists{microtype.sty}{% use microtype if available
  \usepackage[]{microtype}
  \UseMicrotypeSet[protrusion]{basicmath} % disable protrusion for tt fonts
}{}
\makeatletter
\@ifundefined{KOMAClassName}{% if non-KOMA class
  \IfFileExists{parskip.sty}{%
    \usepackage{parskip}
  }{% else
    \setlength{\parindent}{0pt}
    \setlength{\parskip}{6pt plus 2pt minus 1pt}}
}{% if KOMA class
  \KOMAoptions{parskip=half}}
\makeatother
\usepackage{longtable,booktabs,array}
\newcounter{none} % for unnumbered tables
\usepackage{calc} % for calculating minipage widths
% Correct order of tables after \paragraph or \subparagraph
\usepackage{etoolbox}
\makeatletter
\patchcmd\longtable{\par}{\if@noskipsec\mbox{}\fi\par}{}{}
\makeatother
% Allow footnotes in longtable head/foot
\IfFileExists{footnotehyper.sty}{\usepackage{footnotehyper}}{\usepackage{footnote}}
\makesavenoteenv{longtable}
\setlength{\emergencystretch}{3em} % prevent overfull lines
\providecommand{\tightlist}{%
  \setlength{\itemsep}{0pt}\setlength{\parskip}{0pt}}
\usepackage{newunicodechar}
\newfontfamily\numeralfallback{STIX Two Text}
\newunicodechar{∎}{\ensuremath{\blacksquare}}
\usepackage{fvextra}
\DefineVerbatimEnvironment{Highlighting}{Verbatim}{breaklines,breakanywhere,breakanywheresymbolpre={},breakanywheresymbolpost={},breaksymbolleft={},breaksymbolright={},commandchars=\\\{\}}
\newunicodechar{⁰}{{\numeralfallback\char"2070\relax}}
\newunicodechar{₀}{{\numeralfallback\char"2080\relax}}
\newunicodechar{₁}{{\numeralfallback\char"2081\relax}}
\newunicodechar{₂}{{\numeralfallback\char"2082\relax}}
\newunicodechar{₃}{{\numeralfallback\char"2083\relax}}
\newunicodechar{₄}{{\numeralfallback\char"2084\relax}}

\usepackage{seqsplit}
\renewcommand{\texttt}[1]{{\ttfamily\seqsplit{#1}}}
\setlength{\emergencystretch}{6em}

\newunicodechar{⁻}{{\numeralfallback\char"207B\relax}}
\newunicodechar{⁸}{{\numeralfallback\char"2078\relax}}
\newunicodechar{ⁱ}{{\numeralfallback\char"2071\relax}}
\newunicodechar{⁴}{{\numeralfallback\char"2074\relax}}
\newunicodechar{⁷}{{\numeralfallback\char"2077\relax}}
\newunicodechar{⁶}{{\numeralfallback\char"2076\relax}}
\newunicodechar{⁵}{{\numeralfallback\char"2075\relax}}
\newunicodechar{₈}{{\numeralfallback\char"2088\relax}}
\newunicodechar{ᵢ}{\textsubscript{i}}
\newunicodechar{ⱼ}{\textsubscript{j}}
\newunicodechar{ₜ}{\textsubscript{t}}
\newunicodechar{₅}{{\numeralfallback\char"2085\relax}}
\newunicodechar{₆}{{\numeralfallback\char"2086\relax}}
\newunicodechar{₇}{{\numeralfallback\char"2087\relax}}
\newunicodechar{ᵉ}{\textsuperscript{e}}
\newunicodechar{ᶠ}{\textsuperscript{f}}
\usepackage{bookmark}
\IfFileExists{xurl.sty}{\usepackage{xurl}}{} % add URL line breaks if available
\urlstyle{same}
% fallback for those not using the hyperref driver hyperxmp:
\makeatletter
\@ifundefined{xmpquote}{\newcommand{\xmpquote}[1]{#1}}{}
\makeatother
\hypersetup{
  pdftitle={Triangle geometries of finite linear spaces},
  pdfauthor={Human Compression project},
  hidelinks,
  pdfcreator={LaTeX via pandoc}}

\title{Triangle geometries of finite linear spaces}
\author{Human Compression project}
\date{11 October 2026 --- author draft}

\begin{document}
\maketitle

\section{Abstract}\label{abstract}

We recover the principal classification of Delaby, Leemans and Tranchida
{[}DLT{]} through the flag-relation and frame-orbit interface {[}FR{]}.
The geometric construction, its principal classification target, and the
older classification of locally two-transitive linear spaces are
explicitly attributed. The group-action statement requires a correction:
in an affine plane, semilinear frame transitivity is a condition on
determinant orbits, not containment of the full linear group. We give a
uniform necessary and sufficient criterion and a concrete counterexample
over the field of nine elements. This correction does not change the
classified geometries whose triangle systems are firm, residually
connected and flag-transitive. The full automorphism and correlation
groups, the thin cases, and all diagrams in the principal finite family
are determined. The Hermitian connectivity argument reuses an
independently reviewed norm-shell theorem; the finite exceptional
residue is checked directly from its defining design.

\section{1. Exact scope and retained
inputs}\label{exact-scope-and-retained-inputs}

Our source edition is the final journal article {[}DLT{]}, Algebraic
Combinatorics 9(2) (2026), 379--402, not its earlier preprint. Its exact
PDF is bound in the accompanying source contract. The principal targets
are Theorems 1.1 and 1.2, together with the construction, automorphism,
correlation and residue statements needed for them and the diagrams and
Table 1. We correct literal defects rather than treating a false
statement as a theorem we have reproduced. The source's final
experimental questions are a separate track and are not counted as part
of this recovery. The Desargues-configuration and square examples, and
the identification with a numbered hypermap in an external catalogue,
are auxiliary assertions outside this principal contract. We do not
describe this as recovery of every assertion printed in the paper. The
general isomorphism statements between two uniform input spaces are,
however, included in {[}FR, Theorem 2{]}.

We use {[}FR{]} for the full flag relation, reconstruction, correlation
and residue interface, and its determinant-orbit theorem. These are
complete proofs in an actual separate paper. The following retained
results remain substantial mathematical inputs.

\begin{enumerate}
\def\labelenumi{\arabic{enumi}.}
\tightlist
\item
  Buekenhout, Dehaye and Leemans {[}BDL, Theorem 6.2, pp.42--43{]}
  classify finite faithful flag-transitive linear spaces whose point
  stabilizers and line stabilizers act two-transitively on their
  respective pencils and lines. The classification ultimately uses
  finite group classification. We retain its necessity list, including
  its specified natural actions. We check the extra noncollinear-triple
  condition below; we do not retain an unchecked blanket converse. Its
  circle paragraph already assumes three-transitivity.
\item
  The fundamental automorphism theorems for finite projective and affine
  spaces identify their full groups as PΓL and AΓL, respectively. For
  the classical Hermitian unital, the full group is PΓU(3,q), with
  U\_H(2) identified with AG(2,3). These classical identifications are
  retained as in DLT; they are not new group-classification theorems
  proved here.
\item
  The complete Hermitian transport and generation proofs are in
  Part III of the common paper [FR]. Theorems H.1--H.3 prove that a
  nonzero norm shell is a regular SU(2,q)-set and that two specified
  isotropic directions suffice for connectedness. These arguments are
  included in the present topic, rather than imported without their proof.
\item
  The usual Mathieu permutation actions and the classical classification
  of finite three-transitive permutation groups, as retained in BDL's
  circle case, are used with their actual three-transitivity hypothesis.
  In its projective-line branch the required natural-action normalizer
  identity is proved separately in {[}FR, Theorem 6{]}, including q=2,3;
  it is not inferred from the dimension-at-least-two collineation
  theorem. The exceptional A7 linear action needed here is also checked
  directly in the finite certificate.
\end{enumerate}

All remaining implications in the principal contract are proved below or
in {[}FR{]}. A computation for finitely many fields is never substituted
for an assertion for every prime power.

\section{2. Corrected classification of the
actions}\label{corrected-classification-of-the-actions}

For a linear space S, let T(S) be its ordered noncollinear triples. A
circle K\_v means the linear space whose lines are all two-element
subsets of its v-point set. We exclude the degenerate one-line space,
consistently with DLT's requirement that at least two lines pass through
each point.

\textbf{Theorem 1 (corrected action classification).} Let G act
faithfully on a finite linear space S and transitively on T(S). Then
(S,G) belongs to one of the following families, and every specified pair
has this transitivity.

\begin{enumerate}
\def\labelenumi{(\alph{enumi})}
\item
  S=PG(d,q), d≥2, and PSL(d+1,q)≤G≤PΓL(d+1,q).
\item
  S=PG(3,2), with the exceptional natural group G≅A7.
\item
  S=AG(d,q), d≥2 and q≥3, and G=K\^{}d⋊H with SL(d,q)≤H≤ΓL(d,q). For d≥3
  no further restriction is needed. For d=2 the determinant-semilinear
  image of H must be transitive on K*, under t↦det(A)σ(t) for the
  transformation v↦v\^{}σ A.
\item
  S=U\_H(q), q=2 or 4, and G=PΓU(3,q).
\item
  S=K\_v and G is a finite three-transitive permutation group in its
  given v-point action. Its customary list is: S\_v for v≥3; A\_v for
  v≥5; groups between PSL(2,q) and PΓL(2,q) in their projective-line
  action which are three-transitive; the Mathieu actions
  M11,M12,M22,M23,M24 of the indicated degrees and Aut(M22); the
  degree-12 action of M11; AGL(d,2), d≥2; and the exceptional 2\^{}4:A7
  action of degree 16. In the projective-line case every intermediate
  group is three-transitive when q is even. For odd q its projective
  determinant-square image must be nontrivial. Coincidences between
  entries are permitted. There is no A7 example of degree 15.
\end{enumerate}

\emph{Proof of necessity.} Triple transitivity first implies flag
transitivity and the local two-transitivities required by BDL. To see
flag transitivity, extend a flag (p,L) to (p,a,b) with a∈L\{p\} and b
outside L, and use the triple action. For two distinct lines through p,
choose a and b on them other than p.~Sending the resulting triple to any
similarly chosen triple fixes p and sends the ordered pair of lines as
desired. For two distinct points on L, choose a third point outside L;
the same argument gives the two-transitivity of the setwise line
stabilizer on L. Faithfulness is part of our hypothesis, so all the
assumptions of BDL's classification apply.

Besides (a)--(e), BDL has the exceptional pair AG(2,4), AΓL(1,16). Here
there are 16·15·12=2880 ordered noncollinear triples but the acting
group has order 16·15·4=960; it cannot be transitive. In the affine
family BDL provides the full translation subgroup and the contained SL
group. The ambient upper bound is ΓL by the affine automorphism theorem.
The remaining condition is exactly {[}FR, Theorem 5{]}, which proves
(c), including the semilinear plane case. The projective and A7 cases
need no additional necessity restriction. The unital branch is
determined in §3. A circle has no collinear triple of distinct points,
so its condition is precisely three-transitivity. The circle list is the
retained permutation classification, with the explicit impossible
low-degree entries removed. For its PSL-containing natural
projective-line branch, the normalizer identity {[}FR, Theorem 6{]}
supplies the stated PΓL upper bound. This proves necessity, subject only
to the unital calculation supplied next.

\emph{Proof of sufficiency outside the unital case.} In PG(d,q), three
noncollinear points are represented by three independent vectors. Extend
them to a vector basis, and do the same for a target triple. A linear
map between these bases carries the projective triple correctly. Rescale
one target basis vector to make its determinant one; rescaling does not
change that target projective point. Hence PSL acts transitively. For
(b), the certificate realizes GL(4,2) on its 15 nonzero vectors,
identifies its exceptional A7 subgroup, and checks its single orbit of
length 15·14·12=2520 on independent ordered triples. This is the
complete one-field exception, not a sample of an infinite family. For
(c), translating the first point to zero and applying {[}FR, Theorem
5{]} proves sufficiency.

For the circle case, three-transitivity is exactly the desired
assertion. The standard actions in the list are retained as stated,
while the bounds on alternating groups are immediate: S\_v sends every
ordered triple to every other one; when v≥5 its parity can be corrected
on two unused points, so A\_v does as well. The orders of A3 and A4 are
insufficient. On the projective line PGL(2,q) is sharply
three-transitive. For odd q PSL has two orbits on ordered triples,
distinguished by the determinant-square class of the unique projectivity
carrying (∞,0,1) to that triple. A field automorphism fixes the base
triple and preserves square classes. Thus an intermediate semilinear
group is transitive exactly when it has a nonsquare-determinant element.
For even q PSL=PGL. In affine binary space, any two distinct nonzero
vectors are independent; GL(d,2) is transitive on such pairs, proving
the assertion for AGL(d,2). The finite certificate also checks that the
same exceptional A7 in GL(4,2) is two-transitive on the 15 nonzero
vectors, which proves the 2\^{}4:A7 assertion after translations. The
standard Mathieu actions are retained, as declared above. ∎

The impossible A7 circle entry already occurs in the older BDL display
and was already identified in the project's earlier {[}IC, §4{]}. Its
order 2520 is smaller than 15·14·13=2730. This observation does not
affect its valid PG(3,2) action, whose noncollinear count is2520. DLT's
named circle list also permits A\_v for v=3,4; its preliminary
three-transitivity hypothesis excludes those groups, but the displayed
named converse must not be read as asserting their three-transitivity.
These are reused source-audit findings, not newly answered research
questions.

\section{3. The unital orbit
parameter}\label{the-unital-orbit-parameter}

Use homogeneous coordinates over K=F\_(q²) and the Hermitian form

\[
 h(x,x)=x_1x_3^q+x_1^qx_3+x_2^{q+1}.
\]

The projective unitary group is two-transitive on the isotropic points;
this is a standard consequence of Witt extension, retained here. Fix an
ordered pair as e1,e3. Every isotropic point outside their joining line
has a unique representative (x1,1,x3), with both x1 and x3 nonzero, and

\[
 z=x_1x_3^q,\qquad z+z^q=-1.
\tag{9}
\]

The complete PGU stabilizer of e1,e3 is represented by
diag(a,1,a\^{}(−q)), a∈K*. Indeed preservation of the two isotropic
lines also preserves their orthogonal complementary line; after
projective normalization of its middle diagonal entry the Hermitian
condition gives exactly these matrices. They preserve z and act
regularly on each z-fiber: normalize x1 to1 and then x3 is fixed by (9).
In particular there are q PGU orbits on T(S), and the stabilizer of any
ordered noncollinear triple is trivial.

The field group acts on z by its p-power Frobenius, where q=p\^{}e. If q
is odd, z=−1/2 is a fixed member of (9), while the fiber has
q\textgreater1 elements. Thus no subgroup of PΓU is transitive. If
q=2\^{}e with e≥3, every field orbit has length at most2e\textless q,
again excluding transitivity. For q=2 the trace-one fiber consists of
the two elements of F4\textbackslash F2, a single Frobenius orbit. For
q=4 it has four elements, none in F4, since an element of F4 has trace
z+z\^{}4=0. They all have degree four over F2 and form one Frobenius
orbit. The full PΓU group is therefore transitive exactly for q=2,4.

In both cases its order equals \textbar T(S)\textbar:

\[
 |T(S)|=(q^3+1)q^3(q^3-q),\qquad
 |\operatorname{P\Gamma U}(3,q)|=(q^3+1)q^3(q^2-1)\,2e.
\]

Here q=2e for q=2,4. Any transitive subgroup has order at least
\textbar T(S)\textbar, so it must be the full group. This proves both
directions of Theorem 1(d), including the group requirement that
cardinal inequalities alone do not supply for arbitrary q. The small
identification U\_H(2)≅AG(2,3) is also verified by the finite model; it
is not used to infer the q=4 calculation.

\section{4. A semilinear counterexample and the repaired affine
statement}\label{a-semilinear-counterexample-and-the-repaired-affine-statement}

Take K=F9, let z generate K*, and let σ:x↦x³ be its nontrivial field
automorphism. Define H≤ΓL(2,9) by the condition

\[
 (A,\sigma^e)\in H\quad\Longleftrightarrow\quad
 \det A\text{ is a square exactly when }e=0,
 \qquad e\in\{0,1\}.
\tag{10}
\]

This is a subgroup: Frobenius preserves square classes and the
determinant class of a composite is the product of the two determinant
classes. Its linear half contains SL(2,9), of order 720, and has
determinant image the four squares. The other half is obtained using
(diag(z,1),σ). Thus \textbar H\textbar=2·4·720=5760. It omits every
nonsquare-determinant linear matrix and therefore does not contain
GL(2,9), which also has order 5760.

The induced action (8) is transitive on K*: the linear half is
transitive within either square class, and the specified semilinear
element interchanges the two classes. By {[}FR, Theorem 5{]}, H is
transitive on the 5760 ordered bases of K², hence regular there. Adding
all translations gives a group of order81·5760=466560 acting regularly
on the ordered noncollinear triples of AG(2,9). These transformations
preserve every affine line, so this is a counterexample within the exact
hypotheses of DLT Theorem 1.1, not merely an abstract permutation
action. The native GAP certificate verifies the defining line action,
subgroup orders, noncontainment and the entire triple orbit.

The replacement in Theorem 1(c) is uniform for all finite fields and all
semilinear subgroups containing SL. If one adds the extra assumption
that H is linear, it reduces to the original GL conclusion. This
distinction is a source correction and a positive orbit theorem; it is
not an answer to an explicitly open conjecture.

\section{5. Residual connectedness and the classified
geometries}\label{residual-connectedness-and-the-classified-geometries}

\textbf{Theorem 2 (triangle classification).} For a finite linear space
S, Δ(S) is firm, residually connected and flag-transitive if and only if
S is PG(2,q), q≥2; AG(2,q), q≥3; or U\_H(q), q=2 or4. In each of these
cases Δ(S) has a duality exactly when S does.

Suppose first that Δ(S) has the three stated properties. Flag
transitivity forces uniformity by {[}FR, §1{]}, and firmness forces k≥3.
Moreover r≥3: if r=2, v=2k−1 and b=vr/k=4−2/k would be an integer; k≥3
is impossible. Consequently {[}FR, Theorem 2{]} reconstructs the full
automorphism group of S, and its chamber correspondence puts S in
Theorem 1. The affine and projective cases of dimension at least three
have disconnected vertex residues: {[}FR, Theorem 4{]} separates them
into the distinct planes through a fixed line. Circles have k=2 and are
not firm. The remaining list is exactly the one stated.

We prove that all surviving cases have the required properties. The
appropriate groups in Theorem 1 give flag transitivity; k≥3 gives
firmness. Global connectivity follows from {[}FR, §3{]}. It remains to
show that every vertex residue is connected. It suffices to check the
deleted-pencil model {[}FR,(7){]}, either uniformly at every flag or at
one flag when its group is already flag-transitive.

In a projective plane, deleting the line L leaves an affine plane. The
allowed pencils in (7) are all its parallel classes except the class
indexed by p.~Choose two remaining classes. Every line of one meets
every line of the other, so their union and all outside points form a
connected incidence grid. Every additional allowed line meets this grid.
This proves residue connectivity, including order2.

In an affine plane of order q≥3, choose two distinct a,b in L\{p\}.
There are q non-L lines through each. A line through a meets every one
through b except its unique parallel, and those meetings are outside L.
The intersection graph of the two pencils is K\_(q,q) minus a perfect
matching, which is connected for q≥3. Each outside point lies on exactly
one line from each pencil, and each additional allowed line contains
outside points. This proves the full residue connected. The argument
requires no Desarguesian coordinates.

For a Hermitian unital, apply Theorems H.1--H.3 of the common paper.
With L fixed, its outside isotropic points identify with a nonzero norm
shell in a Hermitian two-space, and its points on L identify with
isotropic directions. An edge through a point a of L is the
corresponding isotropic-direction move. Formula (7) deletes precisely
the direction p. There are q+1 isotropic directions, so for q≥2 choose
two distinct allowed directions. Choose representatives e,f of those
lines. Their Hermitian pairing is nonzero, and rescaling gives a
hyperbolic basis with those two axes. Thus the two explicit root
families of [FR, Part III] apply to these particular allowed
directions. They generate SU(2,q); by norm-shell transitivity, the moves
in these allowed directions connect every outside point. Every
allowed-line vertex contains outside points. Thus the entire residue is
connected. This proves the stronger uniform result that Δ(U\_H(q)) is
firm and residually connected for every prime power q≥2, although it is
flag-transitive only for q=2,4. The full transport, root generation and geometric identification are
proved in [FR, Part III]. The allowed directions have been specified
explicitly; the argument does not assume that deleting an arbitrary
part of a connected graph preserves connectedness.

Finally all surviving spaces are thick as rank-two spaces. The duality
assertion is {[}FR, Theorem 3{]}. This completes the proof of Theorem 2.
∎

This proves DLT's geometric principal theorem without changing its list.
In particular AG(2,q) for q\textgreater3 yields a thick triangle
geometry with a triality and no duality: it has q² points and q(q+1)
lines, precluding a duality of its input. U\_H(4) also has unequal point
and block counts65 and208. Projective planes have their usual
polarities.

\section{6. The complete residue
diagrams}\label{the-complete-residue-diagrams}

Write k for the input line size. Each type contains vr flags, and every
panel has k−1 completions. Consequently the vertex labels in the usual
Buekenhout convention are s=k−2 and n=vr. Triality identifies the three
residue types. We give girth, both part-diameters, and all small
exceptions.

For PG(2,q), the residue is the affine plane with one parallel class
removed. When q=2 it is the eight-cycle. For q≥3 its girth is6 and both
part-diameters are4. Indeed any two points either have an allowed
joining line or are joined by a four-step path using two allowed
directions. A point and a nonincident allowed line have distance3, using
another allowed direction through the point. Distinct parallel allowed
lines have distance4 through a line of another class, and nonparallel
lines have distance2. There are at least three allowed directions;
choose a line in each with the third avoiding the intersection of the
first two. Their three pairwise intersections give a six-cycle. No
four-cycle can exist in a linear space. Thus every edge label is
(4,3,4), except that q=2 gives a generalized quadrangle of order(1,1),
with label4. The vertices have n=(q²+q+1)(q+1) and s=q−1.

For AG(2,q), put L=\{y=0\}, p=(0,0). Its residue has outside points
(x,y), y≠0, and lines represented by (a,t), a≠0, t∈K, with incidence

\[
 x=a+ty.
\tag{11}
\]

For q=3, there are two allowed anchors a, and the graph is the
subdivision of K\_(3,3) minus a perfect matching: a twelve-cycle. For
q≥4, every two point vertices have distance at most4. To see this when
there is no common allowed line, choose a line through the first point
with anchor a. Among the q−2 allowed anchors b≠a for lines through the
second point, at most one yields a parallel line. At least one other
choice remains, and its intersection with the first line is outside L.
This supplies the desired four-step path. For two line vertices lacking
an outside intersection, choose an outside point Q on the first. Of the
q−1 outside points on the second, at most one has the same y-coordinate
as Q and at most one joins Q to p.~Any other choice gives an allowed
joining line and a four-step path. Similarly a point and a nonincident
allowed line have distance3: choose an anchor different from the line's
anchor and avoid the at most one parallel choice. These counts hold
already for q=4.

There exist point pairs on a horizontal line whose joining line is not
allowed, and two distinct residue lines through the same allowed anchor;
both exhibit distance4. There are no four-cycles. To obtain a six-cycle,
choose two nonparallel allowed lines with distinct anchors, then a third
anchor distinct from both. Of the q lines at that anchor, at most two
are parallel to the first pair and at most one passes through their
outside intersection. Since q≥4 a remaining line forms a triangle wholly
outside L. Thus the girth is6 and both part-diameters are4. Every edge
label is (4,3,4), except for q=3 where it is the generalized-hexagon
label6. The vertices have n=q²(q+1) and s=q−2.

For U\_H(2), use its identification with AG(2,3). For U\_H(4), construct
the 65 isotropic points in PG(2,16) and its208 secant lines directly.
Delete one pencil exactly as in (7). The resulting graph has60 vertices
in each part, degree 4, and240 edges. The certificate performs
breadth-first search from every vertex, finding the following complete
distance profiles:

{\def\LTcaptype{none} % do not increment counter
\begin{longtable}[]{@{}ll@{}}
\toprule\noalign{}
Part & Numbers of vertices at distances 0,1,2,3,4,5,6 \\
\midrule\noalign{}
\endhead
\bottomrule\noalign{}
\endlastfoot
outside-point vertices & 1,4,12,36,44,20,3 \\
allowed-line vertices & 1,4,12,36,42,20,5 \\
\end{longtable}
}

The same searches give shortest circuit length 8. Hence both diameters
are 6 and the gonality is 4. Moreover these two different profiles are
automorphism invariants, so no graph automorphism interchanges the
parts: the rank-two residue has no duality. The full group is
flag-transitive by §3, so the calculation at one flag gives every vertex
residue, and triality gives all type pairs. The three diagram edges are
(6,4,6), and the vertex labels are n=65·16=1040 and s=3. This is a
complete bounded certificate for the only exceptional diagram, with no
group-order guess or sampled-vertex inference.

\section{7. Full groups, higher dimensions and the thin
cases}\label{full-groups-higher-dimensions-and-the-thin-cases}

For all uniform thick cases the exact correlation description is
C3⋊Cor(S), with the sign action in {[}FR, Theorem 3{]}. The classical
full automorphism identifications in §1 therefore give the following
table. The symbol ``thin'' refers to Δ, not to its rank-two input.

{\small\def\LTcaptype{none} % do not increment counter
\begin{longtable}[]{@{}
  >{\raggedright\arraybackslash}p{(\linewidth - 10\tabcolsep) * \real{0.1429}}
  >{\centering\arraybackslash}p{(\linewidth - 10\tabcolsep) * \real{0.1905}}
  >{\centering\arraybackslash}p{(\linewidth - 10\tabcolsep) * \real{0.1905}}
  >{\centering\arraybackslash}p{(\linewidth - 10\tabcolsep) * \real{0.1905}}
  >{\raggedright\arraybackslash}p{(\linewidth - 10\tabcolsep) * \real{0.1429}}
  >{\raggedright\arraybackslash}p{(\linewidth - 10\tabcolsep) * \real{0.1429}}@{}}
\toprule\noalign{}
\begin{minipage}[b]{\linewidth}\raggedright
Input S
\end{minipage} & \begin{minipage}[b]{\linewidth}\centering
Connected
\end{minipage} & \begin{minipage}[b]{\linewidth}\centering
Residually connected
\end{minipage} & \begin{minipage}[b]{\linewidth}\centering
Thin
\end{minipage} & \begin{minipage}[b]{\linewidth}\raggedright
Aut(Δ)
\end{minipage} & \begin{minipage}[b]{\linewidth}\raggedright
Cor(Δ)
\end{minipage} \\
\midrule\noalign{}
\endhead
\bottomrule\noalign{}
\endlastfoot
AG(d,q), d≥2,q≥3 & yes & exactly d=2 & exactly q=3 & AΓL(d,q) &
AΓL(d,q)×C3 \\
PG(d,q), d≥2 & yes & exactly d=2 & exactly q=2 & PΓL(d+1,q) & see
below \\
U\_H(4) & yes & yes & no & PΓU(3,4) & PΓU(3,4)×C3 \\
K\_v, v≥4 & yes & no & no & S\_v & S\_v×S3 \\
K\_3 & no & no & no & S6 & S6×S3 \\
\end{longtable}
}

For PG(2,q), its standard polarity gives Cor(Δ)≅PΓL(3,q)⋊S3; here C3
centralizes the collineation group and the transposition acts by that
polarity. For d≥3 there are different numbers of points and lines, so
Cor(Δ)=PΓL(d+1,q)×C3. This also supplies the higher-dimensional q=2 rows
that are not all printed in DLT's table. The U\_H(2) row duplicates
AG(2,3). The affine rows require AΓL, including field automorphisms. For
a direct witness when q is not prime, Frobenius preserves every affine
line but cannot be K-affine linear: it fixes 0 and the standard basis,
while moving z e1 for a field element z outside the prime field. In
particular the literal AGL label in DLT Table 1 omits actual
automorphisms.

For completeness we prove all circle statements. Theorem 2 of {[}FR{]}
applies when v≥4, giving Aut(Δ(K\_v))=S\_v. The map

\[
 (p,\{p,q\},i)\longmapsto(q,\{p,q\},-i)
\tag{12}
\]

is an incidence-preserving involution: it reverses the endpoint order of
each directed edge and reverses the cyclic type order. Together with τ
it is a copy of S3, commuting with the diagonal S\_v. Every possible
type permutation is represented and the kernel is exactly S\_v, proving
Cor(Δ)=S\_v×S3.

Its underlying directed flag relation is the nonbacktracking arc
relation on K\_v: (q,p) advances to (p,a) with a distinct from q,p.~For
v≥4 its arc graph is strongly connected. For example an arc can be
reversed by the walk q,p,a,b,p,q, using two additional distinct
vertices, after which the ordinary connectedness of K\_v permits
movement to any arc without an immediate reversal. Based at any arc
there are directed cycles of lengths three and four. Lifting to the
cyclic type coordinate, the four-cycle changes the type by one, while
returning to the same arc. Thus the lift Δ is connected. In its vertex
residue, (7) reduces to v−2 disjoint edges: an outside point Q is
incident only to the allowed line qQ, where L=\{p,q\}. It is not
residually connected and its panels have only one completion, so it is
neither firm nor thin. When v=3, every flag extends to exactly one
chamber and Δ is the disjoint union of its six chambers. Its
type-preserving automorphisms permute those chambers arbitrarily, giving
S6; the additional type permutations give the direct product S6×S3. This
proves the entire small-case table without applying a thick
reconstruction theorem outside its hypotheses.

\section{8. Literal source corrections and proof
coverage}\label{literal-source-corrections-and-proof-coverage}

The principal geometry classification, construction, triality and source
authorship remain unchanged. The corrected action classification
replaces one affine-plane necessity by the exact determinant criterion.
The finite circle-list defects were already recorded in {[}IC{]}; no new
question credit is attached to them. The affine full-group table
includes field automorphisms. The general correlation formula C3⋊Cor(S)
does not require a polarity, whereas splitting its natural quotient onto
S3 does.

There is one further logical limitation in DLT's description of its
construction as a functor. Its §2 morphisms need preserve incidence only
in the forward direction, while Construction 1 uses both incidence and
nonincidence. Such a morphism does not in general induce a triangle
morphism. Indeed map every input point to a fixed point p and every line
to a fixed incident line L. This is a type-preserving incidence
morphism, but its proposed triangle map sends each type to (p,L,i), and
no such vertices of different types are incident. This loses every
source edge. The construction does transport correlations, as proved in
{[}FR{]}, and more generally maps preserving both incidence and
nonincidence. No use of the false general functor assertion enters
either principal proof.

The Lean file proves actual equal-determinant matrix transport,
semilinear determinant transformation and subgroup orbit saturation. It
does not formalize BDL's classification, the finite geometry model, all
residue connectivity, or all source theorem instances. The finite
certificate proves the complete fixed exceptional models and advertised
controls, not the all-field classification on its own. The complete
written proof consists of this application, its full {[}FR{]} dependency
and the expressly retained inputs. A future matched compression
measurement must charge all of them under the agreed topic accounting
rule. No ratio, principal-source admission or open-question count is
inferred from the number of lemmas or from a successful executable
replay.

\section{References}\label{references}

{[}DLT{]} R. Delaby, D. Leemans and P. Tranchida, \emph{Geometries with
trialities arising from linear spaces}, Algebraic Combinatorics 9(2)
(2026),379--402, DOI10.5802/alco.482.

{[}BDL{]} F. Buekenhout, P.-O. Dehaye and D. Leemans, \emph{RWPRI and
(2T)1 flag-transitive linear spaces}, Beiträge zur Algebra und
Geometrie44(1) (2003),25--46. Theorem 6.2, printed42--43, is a retained
classification.

{[}FR{]} \emph{Flag relations and derived incidence geometries},
the common paper of this topic. Numbered results 1--6 are in Part I;
results H.1--H.3 are in Part III.

{[}IC{]} \emph{Incidence compression}, earlier repository manuscript,
§4. Its matrix reconstruction, correlation extension, residue bridge and
already recorded source corrections are reused with attribution.

{[}HU{]} The project's earlier proof for Hermitian unital residues and
its independent review. Its complete transport and generation arguments
are reproduced in [FR, Part III], with the geometric identification for
the triangle construction.

\end{document}
